% ===================================================================== PREAMBULE COMMUN — Carnet d'entraînement Fiche 9
\documentclass[11pt,a4paper]{article}
\usepackage[utf8]{inputenc}
\usepackage[T1]{fontenc}
\usepackage{lmodern}
\usepackage[french]{babel}
\usepackage{amsmath,amssymb,mathtools}
\usepackage{icomma}
\usepackage{siunitx}
\sisetup{output-decimal-marker={,},per-mode=symbol}
\DeclareSIUnit{\molar}{M}
\DeclareSIUnit{\Molar}{mol\,L^{-1}}
\usepackage[version=4]{mhchem}
\usepackage{xcolor}
\usepackage{colortbl}
\usepackage{tikz}
\usepackage{pgfplots}
\usepackage[most]{tcolorbox}
\usepackage{enumitem}
\usepackage{array}
\usepackage{booktabs}
\usepackage{multicol}
\usepackage{fancyhdr}
\usepackage[hidelinks]{hyperref}
\usetikzlibrary{arrows.meta,positioning,calc,decorations.pathmorphing,
  decorations.markings,angles,quotes,patterns,backgrounds,
  shapes.geometric,shapes.misc,fadings,intersections,fit}
\pgfplotsset{compat=1.18}
\usepackage[a4paper,margin=2.2cm,headheight=26pt,headsep=10pt]{geometry}

% ---------- Palette ----------
\definecolor{primary}{RGB}{150,98,162}
\definecolor{primarydark}{RGB}{92,52,110}
\definecolor{defcol}{RGB}{0,114,178}
\definecolor{thmcol}{RGB}{0,140,120}
\definecolor{formulacol}{RGB}{198,93,9}
\definecolor{remarkcol}{RGB}{120,94,200}
\definecolor{attncol}{RGB}{200,40,55}
\definecolor{excol}{RGB}{0,135,90}
\definecolor{methcol}{RGB}{90,90,100}
\definecolor{cationcol}{RGB}{0,90,190}
\definecolor{anioncol}{RGB}{200,60,55}
\definecolor{cristalcol}{RGB}{110,105,120}
\definecolor{eaucol}{RGB}{120,180,220}
\definecolor{solcol}{RGB}{0,135,90}
\definecolor{kpscol}{RGB}{198,93,9}
\definecolor{qcol}{RGB}{120,94,200}
\definecolor{precipcol}{RGB}{110,70,40}
\definecolor{exercol}{RGB}{150,98,162}
\definecolor{corrcol}{RGB}{0,120,100}

% ---------- Macros ----------
\newcommand{\Kps}{K_{\mathrm{ps}}}
\newcommand{\Ce}[1]{\left[\,\ce{#1}\,\right]}

% ---------- Style des boîtes ----------
\tcbset{boxbase/.style={enhanced,breakable,boxrule=0.9pt,arc=2.5pt,
   left=3mm,right=3mm,top=2.3mm,bottom=2.3mm,
   fonttitle=\bfseries,coltitle=white,toptitle=0.6mm,bottomtitle=0.6mm}}

\newtcolorbox[auto counter]{definition}[1][]{%
  boxbase,colback=defcol!5,colframe=defcol,
  title={\textsc{Définition}~\thetcbcounter\ \ifx\relax#1\relax\relax\else\textmd{--- \itshape #1}\fi}}
\newtcolorbox[auto counter]{theoreme}[1][]{%
  boxbase,colback=thmcol!6,colframe=thmcol,
  title={\textsc{Loi / Propriété}~\thetcbcounter\ \ifx\relax#1\relax\relax\else\textmd{--- \itshape #1}\fi}}
\newtcolorbox[auto counter]{formule}[1][]{%
  boxbase,colback=formulacol!6,colframe=formulacol,
  title={\textsc{Formule}~\thetcbcounter\ \ifx\relax#1\relax\relax\else\textmd{--- \itshape #1}\fi}}
\newtcolorbox{remarque}[1][]{%
  boxbase,colback=remarkcol!6,colframe=remarkcol,
  title={\textsc{Remarque}\ifx\relax#1\relax\relax\else\textmd{ : \itshape #1}\fi}}
\newtcolorbox{attention}[1][]{%
  boxbase,colback=attncol!6,colframe=attncol,
  title={\textsc{Attention}\ifx\relax#1\relax\relax\else\textmd{ : \itshape #1}\fi}}
\newtcolorbox{exemple}[1][]{%
  boxbase,colback=excol!5,colframe=excol,
  title={\textsc{Exemple}\ifx\relax#1\relax\relax\else\textmd{ : \itshape #1}\fi}}
\newtcolorbox{methode}[1][]{%
  boxbase,colback=methcol!7,colframe=methcol,sharp corners,
  title={\textsc{Méthode}\ifx\relax#1\relax\relax\else\textmd{ : \itshape #1}\fi}}
\newtcolorbox{astuce}[1][]{%
  boxbase,colback=primary!7,colframe=primary,
  title={\textsc{Astuce}\ifx\relax#1\relax\relax\else\textmd{ : \itshape #1}\fi}}

\newtcolorbox[auto counter]{exercice}[1][]{%
  boxbase,colback=exercol!4,colframe=exercol,
  title={\textsc{Exercice}~\thetcbcounter\ \ifx\relax#1\relax\relax\else\textmd{--- \itshape #1}\fi}}
\newtcolorbox[auto counter]{correction}[1][]{%
  boxbase,colback=corrcol!4,colframe=corrcol,
  title={\textsc{Corrigé}~\thetcbcounter\ \ifx\relax#1\relax\relax\else\textmd{--- \itshape #1}\fi}}

% ---------- Environnement « où : » ----------
\newenvironment{ou}{%
  \par\smallskip\noindent\textbf{où :}\nopagebreak
  \begin{itemize}[leftmargin=1.8em,topsep=2pt,itemsep=1.5pt,parsep=0pt]}%
  {\end{itemize}}

% ---------- Styles TikZ ----------
\tikzset{
  cat/.style={circle,fill=cationcol,draw=cationcol!50!black,inner sep=0pt,
              minimum size=5mm,font=\bfseries\scriptsize,text=white},
  an/.style={circle,fill=anioncol,draw=anioncol!50!black,inner sep=0pt,
             minimum size=5mm,font=\bfseries\scriptsize,text=white},
  crx/.style={circle,fill=cristalcol,draw=black!50,inner sep=0pt,
              minimum size=5mm,font=\bfseries\scriptsize,text=white},
  ptn/.style={circle,fill=black,inner sep=1.1pt}}

% ---------- En-têtes ----------
\pagestyle{fancy}
\fancyhf{}
\fancyhead[L]{\small\textcolor{primarydark}{\textbf{Fiche 9} — Solubilité et produit de solubilité}}
\fancyhead[R]{\small\textcolor{primarydark}{\textsc{Carnet d'entraînement}}}
\fancyfoot[C]{\small\thepage}
\renewcommand{\headrulewidth}{0.8pt}
\renewcommand{\headrule}{\hbox to\headwidth{\color{primary}\leaders\hrule height \headrulewidth\hfill}}
\usepackage{titlesec}
\titleformat{\section}{\Large\bfseries\color{primarydark}}{\thesection}{0.6em}{}
\titleformat{\subsection}{\large\bfseries\color{primary}}{\thesubsection}{0.6em}{}

% ---------- Bandeau de titre ----------
% #1 = thème/étiquette   #2 = titre principal   #3 = sous-titre
\newcommand{\bandeau}[3]{%
\begin{center}
\begin{tikzpicture}
\node[rectangle,rounded corners=7pt,fill=primary,text=white,
      minimum width=15.6cm,minimum height=1.95cm,align=center,inner sep=8pt]
  {{\large\bfseries #1}\\[3pt]{\Large\bfseries #2}\\[3pt]{\small #3}};
\end{tikzpicture}
\end{center}
\vspace{2mm}}
\begin{document}
\bandeau{Thème 5 — Influence du pH}%
{Exercices — Niveau Difficile}%
{Séparation sélective, prédiction d'acidification, pont ion commun $\leftrightarrow$ pH}

\noindent\textit{Consigne : sous-questions progressives. Données à \qty{25}{\degreeCelsius},
$K_e=\num{1{,}0e-14}$.}
\medskip

\begin{exercice}[séparer le fer du magnésium par le pH]
Une solution contient $[\ce{Fe^3+}]=[\ce{Mg^2+}]=\qty{1{,}0e-2}{\molar}$. On veut précipiter
\emph{sélectivement} le fer. Données : $\Kps(\ce{Fe(OH)3})=\num{2{,}0e-39}$,
$\Kps(\ce{Mg(OH)2})=\num{9{,}0e-12}$. La figure donne la solubilité $s$ (concentration du cation
restant en solution) en fonction du pH.

\begin{center}
\begin{tikzpicture}
\begin{axis}[
   width=11cm,height=5.4cm,
   xlabel={pH}, ylabel={$s$ (\si{\molar})},
   xmin=2,xmax=14,ymin=1e-14,ymax=1,ymode=log,
   axis lines=left, xtick={2,4,6,7,8,10,12,14},
   grid=both, grid style={black!10},
   legend pos=north east, legend style={font=\scriptsize,fill=white}]
  \addplot[thick,anioncol] coordinates {(2,2e-3) (3,2e-6) (4,2e-9) (5,2e-12) (5.5,6.3e-14)};
  \addlegendentry{\ce{Fe(OH)3}}
  \addplot[thick,solcol] coordinates {(8.5,0.9) (9,9e-2) (10,9e-4) (11,9e-6) (12,9e-8) (13,9e-10) (14,9e-12)};
  \addlegendentry{\ce{Mg(OH)2}}
\end{axis}
\end{tikzpicture}\\[-1mm]
{\scriptsize\itshape Repère pH $=7$ (centre) : \ce{Fe(OH)3} est tout en bas (précipité), \ce{Mg(OH)2} hors cadre en haut (encore dissous).}
\end{center}

\begin{enumerate}[label=\textbf{\alph*)},leftmargin=2.2em,topsep=2pt,itemsep=2pt]
  \item Calcule le pH de début de précipitation de \ce{Fe(OH)3}, puis celui de \ce{Mg(OH)2}.
  \item À l'aide de la figure, identifie l'intervalle de pH où le fer est (quasi) totalement précipité
        alors que le magnésium reste dissous.
  \item On fixe le pH à \textbf{7}. Calcule la concentration résiduelle $[\ce{Fe^3+}]$ restant en solution.
  \item Vérifie qu'à pH $=7$ le magnésium ne précipite pas (calcule $Q$ pour \ce{Mg(OH)2} et compare à
        son $\Kps$). Conclus sur la faisabilité de la séparation.
\end{enumerate}
\end{exercice}

\begin{exercice}[prédire l'effet d'une acidification]
On acidifie progressivement (on \emph{baisse} le pH) des solutions saturées des sels suivants :
\[ \ce{AgCl},\quad \ce{CaCO3},\quad \ce{Fe(OH)3},\quad \ce{BaSO4},\quad \ce{FeS}. \]
\begin{enumerate}[label=\textbf{\alph*)},leftmargin=2.2em,topsep=2pt,itemsep=2pt]
  \item Pour chaque sel, dis si la solubilité \textbf{augmente}, \textbf{diminue} ou \textbf{ne change
        pas}, et classe le sel dans l'un des trois cas (hydroxyde / sel d'acide faible / indifférent).
  \item Pour \ce{Fe(OH)3} ($\Kps=\num{2{,}0e-39}$), la solubilité vaut $s=\Kps/[\ce{OH-}]^{3}$. Par quel
        facteur $s$ est-elle multipliée quand on passe de pH $=7$ à pH $=2$ ?
\end{enumerate}
\end{exercice}

\begin{exercice}[deux visions d'un même geste : ion commun \emph{ou} pH]
On part d'une solution où $[\ce{Mg^2+}]=\qty{1{,}0e-3}{\molar}$. On ajoute de la soude \ce{NaOH} jusqu'à
$[\ce{OH-}]=\qty{1{,}0e-3}{\molar}$. On donne $\Kps(\ce{Mg(OH)2})=\num{9{,}0e-12}$.
\begin{enumerate}[label=\textbf{\alph*)},leftmargin=2.2em,topsep=2pt,itemsep=2pt]
  \item Calcule le quotient $Q=[\ce{Mg^2+}][\ce{OH-}]^{2}$ et compare-le à $\Kps$ : y a-t-il précipitation ?
  \item Exprime la situation en termes de \textbf{pH} : quel est le pH de la solution ?
  \item Ajouter des \ce{OH-}, est-ce un \emph{effet d'ion commun} (Thème 4) ou un \emph{effet de pH}
        (Thème 5) ? Explique pourquoi ce sont \textbf{deux descriptions du même phénomène} pour un
        hydroxyde.
\end{enumerate}
\end{exercice}

\end{document}
